\Needspace{28\baselineskip}
\SixthArticle{48}

Es wird sich lohnen, zu zeigen, wie mithilfe der Umformungen 82,91 alle in Artikel 5 gesammelten Formeln sehr leicht allein aus dem binomischen Lehrsatz abgeleitet werden können. Die Formeln I---IV folgen daraus unmittelbar. Die Formeln VI---IX folgen von selbst, wenn \(e^x\) als Grenzwert der Potenz \(\left(1+\frac{x}{i}\right)^i\) oder \(\left(1-\frac{x}{i}\right)^{-i}\) und \(\log x\) als Grenzwert von \(i(x^{\frac{1}{i}}-1)\) betrachtet wird, während \(i\) ins Unendliche wächst. Ferner folgen aus
\[
\begin{aligned}[c]
\cos n\varphi+\BareRoot{-1}.\sin n\varphi&=(\cos\varphi+\BareRoot{-1}.\sin\varphi)^n\\
\cos n\varphi-\BareRoot{-1}.\sin n\varphi&=(\cos\varphi-\BareRoot{-1}.\sin\varphi)^n
\end{aligned}
\]
durch Subtraktion und Addition die Formeln XVIII und XXII und daraus nach Formel 82 sogleich XIX und XXIII; daraus wiederum werden nach Formel 91 XVI, XVII, XX und XXI abgeleitet. Setzt man \(t\) anstelle von \(nt\) und nimmt \(n\) unendlich an, so folgen aus XVI und XX die Formeln XI und XII; nimmt man dagegen \(n\) unendlich klein an, so folgen aus XVI---XVIII die Formeln XIII---XV.

\SixthArticle{49}

Aus der Substitution \(x=\frac{1}{y}\) ergibt sich ebenso
\[
\begin{aligned}&0=\alpha\beta P-(\alpha+\beta-1-(\gamma-2)y)y\frac{\mathrm{d}P}{\mathrm{d}y}\\&+(yy-y^3)\frac{\mathrm{d}\mathrm{d}P}{\mathrm{d}y^2}\end{aligned}
\]
Setzt man dann \(P=y^\mu P'\), so ergibt sich
\[
\begin{aligned}[c]
&\mathrm{I.}\qquad 0=P'(\alpha\beta-\mu(\alpha+\beta-1)\\&\qquad+\mu(\gamma-2)y+(\mu\mu-\mu)(1-y))\\
&-\frac{\mathrm{d}P'}{\mathrm{d}y}(\alpha+\beta-1-(\gamma-2)y-2\mu(1-y))y\\
&+(yy-y^3)\frac{\mathrm{d}\mathrm{d}P'}{\mathrm{d}y^2}
\end{aligned}
\]
Damit der Faktor von \(P'\) durch \(y\) teilbar wird, muss man entweder \(\mu=\alpha\) oder \(\mu=\beta\) setzen; der erste Wert liefert
\[
\begin{aligned}&\mathrm{II.}\qquad
0=P'\alpha(\gamma-\alpha-1)-\frac{\mathrm{d}P'}{\mathrm{d}y}(\beta-\alpha-1-(\gamma-2\alpha-2)y)\\&+(y-yy)\frac{\mathrm{d}\mathrm{d}P'}{\mathrm{d}y^2}\end{aligned}
\]
deren \textit{partikuläre} Lösung lautet
\[
P'=F(\alpha,\alpha+1-\gamma,\alpha+1-\beta,y)
\]
\par\SourceLeaf{229}
Die Gleichung I wird also durch die partikuläre Lösung
\[
P=y^\alpha F(\alpha,\alpha+1-\gamma,\alpha+1-\beta,y).
\]
erfüllt, und ebenso liefert der andere Wert \(\mu=\beta\) eine andere partikuläre Lösung
\[
P=y^\beta F(\beta,\beta+1-\gamma,\beta+1-\alpha,y)
\]
woraus sich die allgemeine Lösung ergibt:
\[
\begin{aligned}&P=Ay^\alpha F(\alpha,\alpha+1-\gamma,\alpha+1-\beta,y)\\&+By^\beta F(\beta,\beta+1-\gamma,\beta+1-\alpha,y)\end{aligned}
\]
Hier bezeichnen \(A,B\) Konstanten, die jedoch nicht beliebig, sondern vollständig bestimmt sind, da \(P\) nicht die allgemeine Lösung der Gleichung 80, sondern nur eine partikuläre ist. Damit uns aber die Bestimmung der Werte der Konstanten \(A,B\) nicht auf unnütze Umwege führt, werden wir dieselbe Gleichung auf andere Weise aus dem bereits Entwickelten ableiten.

Setzt man in Gleichung 91 \(-\frac{x}{1-x}=1-z\) und ersetzt in Gleichung 86 \(\beta\) durch \(\gamma-\beta\), \(\gamma\) durch \(\alpha+1-\beta\) und \(x\) durch \(z\), so folgt
\[
\begin{aligned}[c]
&(1-x)^\alpha F(\alpha,\beta,\gamma,x)\\&=
\frac{\Pi(\gamma-1)\Pi(\beta-\alpha-1)}{\Pi(\gamma-\alpha-1)\Pi(\beta-1)}F(\alpha,\gamma-\beta,\alpha+1-\beta,z)\\
&+\frac{\Pi(\gamma-1)\Pi(\alpha-\beta-1)}{\Pi(\alpha-1)\Pi(\gamma-\beta-1)}z^{\beta-\alpha}F(\beta,\gamma-\alpha,\beta+1-\alpha,z)
\end{aligned}
\]
Nach den Gleichungen 91, 92 gilt aber
\[
\begin{aligned}[c]
&F(\alpha,\gamma-\beta,\alpha+1-\beta,z)\\&=(1-z)^{-\alpha}F\left(\alpha,\alpha+1-\gamma,\alpha+1-\beta,-\frac{z}{1-z}\right)\\[5pt]
&F(\beta,\gamma-\alpha,\beta+1-\alpha,z)\\&=(1-z)^{-\beta}F\left(\beta,\beta+1-\gamma,\beta+1-\alpha,-\frac{z}{1-z}\right)
\end{aligned}
\]
Setzen wir diese Ausdrücke sowie \(z=\frac{1}{1-x}\), \(1-z=-\frac{x}{1-x}\), \(-\frac{z}{1-z}=\frac{1}{x}\) ein, so haben wir
\[
\begin{aligned}[c]
&[93]\qquad F(\alpha,\beta,\gamma,x)\\&=
\frac{\Pi(\gamma-1)\Pi(\beta-\alpha-1)}{\Pi(\gamma-\alpha-1)\Pi(\beta-1)}(-x)^{-\alpha}F\left(\alpha,\alpha+1-\gamma,\alpha+1-\beta,\frac{1}{x}\right)\\
&+\frac{\Pi(\gamma-1)\Pi(\alpha-\beta-1)}{\Pi(\alpha-1)\Pi(\gamma-\beta-1)}(-x)^{-\beta}F\left(\beta,\beta+1-\gamma,\beta+1-\alpha,\frac{1}{x}\right)
\end{aligned}
\]
\Needspace{12\baselineskip}
Dies stimmt mit der oben gefundenen Gleichung überein, wenn man
\[
\begin{aligned}[c]
A&=\frac{\Pi(\gamma-1)\Pi(\beta-\alpha-1)}{\Pi(\gamma-\alpha-1)\Pi(\beta-1)}(-1)^\alpha\\
B&=\frac{\Pi(\gamma-1)\Pi(\alpha-\beta-1)}{\Pi(\alpha-1)\Pi(\gamma-\beta-1)}(-1)^\beta
\end{aligned}
\]
setzt.
\SourceLeaf{230}
Dabei ist zu beachten, dass
\[
\begin{array}{rcl}
(-1)^\alpha&=&\cos\alpha k\pi+\BareRoot{-1}\,.\sin\alpha k\pi\\
(-1)^\beta&=&\cos\beta k\pi+\BareRoot{-1}\,.\sin\beta k\pi
\end{array}
\]
gilt, wobei \(k\) ohne nähere Festlegung eine beliebige ungerade ganze Zahl bezeichnet.

\SixthArticle{50}
Durch Gleichung 93 wird der Wert unserer Funktion für Werte des vierten Elements, die größer als eins sind, auf den Fall zurückgeführt, in dem das vierte Element kleiner als eins ist. Zugleich ist ersichtlich, dass zu \textit{negativen} Werten des vierten Elements, die größer als eins sind, stets ein reeller Wert der Funktion \(F\) gehört, während zu positiven Werten nur dann ein reeller Funktionswert gehören kann, wenn \(\alpha\) und \(\beta\) entweder ganze Zahlen oder rationale Brüche mit ungeraden Nennern sind; in den übrigen Fällen lässt \(F(\alpha,\beta,\gamma,x)\) für einen positiven Wert von \(x\) größer als eins nur imaginäre Werte zu.

\SixthArticle{51}
Die bisher entwickelten Beziehungen zwischen mehreren Funktionen \(F\) waren sämtlich linear; wir fügen eine andere von verschiedener Art hinzu. Es sei
\[
\begin{array}{rcl}
P&=&F(\alpha,\beta,\gamma,x)\\
Q&=&x^{1-\gamma}F(\alpha+1-\gamma,\beta+1-\gamma,2-\gamma,x)\\
R&=&F(\alpha,\beta,\alpha+\beta+1-\gamma,1-x)
\end{array}
\]
sodass \(P,\ Q.\ R\) drei partikuläre Lösungen der Gleichung 80 sind, oder
\[
\begin{aligned}
\mathrm{I.}\quad&0=\alpha\beta P-(\gamma-(\alpha+\beta+1)x)\dfrac{\mathrm dP}{\mathrm dx}\\&\qquad-(x-xx)\dfrac{\mathrm{dd}P}{\mathrm dx^2}\\[4pt]
\mathrm{II.}\quad&0=\alpha\beta Q-(\gamma-(\alpha+\beta+1)x)\dfrac{\mathrm dQ}{\mathrm dx}\\&\qquad-(x-xx)\dfrac{\mathrm{dd}Q}{\mathrm dx^2}\\[4pt]
\mathrm{III.}\quad&0=\alpha\beta R-(\gamma-(\alpha+\beta+1)x)\dfrac{\mathrm dR}{\mathrm dx}\\&\qquad-(x-xx)\dfrac{\mathrm{dd}R}{\mathrm dx^2}
\end{aligned}
\]
Multipliziert man die erste Gleichung mit \(Q\) und die zweite mit \(P\) und subtrahiert, so ergibt sich
\[
\begin{aligned}&0=(\gamma-(\alpha+\beta+1)x)\frac{Q\mathrm dP-P\mathrm dQ}{\mathrm dx}\\&+(x-xx)\frac{Q\mathrm{dd}P-P\mathrm{dd}Q}{\mathrm dx^2}\end{aligned}
\]
Diese Gleichung wird durch Multiplikation mit \(x^{\gamma-1}(1-x)^{\alpha+\beta-\gamma}\) integrierbar und liefert
\SourceLeaf{231}
\[
[94]\qquad A=x^\gamma(1-x)^{\alpha+\beta+1-\gamma}\frac{Q\mathrm dP-P\mathrm dQ}{\mathrm dx}
\]
Ganz ebenso ergibt sich
\[
\begin{array}{rl}
[95]&B=x^\gamma(1-x)^{\alpha+\beta+1-\gamma}\dfrac{R\mathrm dQ-Q\mathrm dR}{\mathrm dx}\\[11pt]
[96]&C=x^\gamma(1-x)^{\alpha+\beta+1-\gamma}\dfrac{R\mathrm dP-P\mathrm dR}{\mathrm dx}
\end{array}
\]
Die Konstanten \(A,\ B,\ C\) lassen sich nach der folgenden Methode leicht bestimmen.

Für \(x=0\) gilt \(P=1\); ferner wird \(x^\gamma Q=xF(\alpha+1-\gamma,\beta+1-\gamma,2-\gamma,x)\) zu \(=0\) für \(x=0\); sein durch \(\mathrm dx\) geteiltes Differential, nämlich \(\gamma x^{\gamma-1}Q+x^\gamma\frac{\mathrm dQ}{\mathrm dx}\), wird jedoch \(=1\). Daraus folgt \(\frac{x^\gamma\mathrm dQ}{\mathrm dx}=1-\gamma\) für \(x=0\) und somit
\[
A=\gamma-1
\]
Um auch \(B\) und \(C\) zu bestimmen, nehmen wir die folgende Gleichung wieder auf:
\[
R=f(\alpha,\beta,\gamma)P+f(\alpha+1-\gamma,\beta+1-\gamma,2-\gamma)Q
\]
Durch Differenzieren erhält man daraus
\[
\frac{\mathrm dR}{\mathrm dx}=f(\alpha,\beta,\gamma)\frac{\mathrm dP}{\mathrm dx}+f(\alpha+1-\gamma,\beta+1-\gamma,2-\gamma)\frac{\mathrm dQ}{\mathrm dx}
\]
Multipliziert man die erste mit \(\frac{\mathrm dQ}{\mathrm dx}\) und die zweite mit \(Q\) und subtrahiert, so ergibt sich
\[
\frac{Q\mathrm dR-R\mathrm dQ}{\mathrm dx}=f(\alpha,\beta,\gamma)\frac{Q\mathrm dP-P\mathrm dQ}{\mathrm dx}\quad\mbox{somit}
\]
\[
B=(1-\gamma)f(\alpha,\beta,\gamma)=\frac{\Pi(\alpha+\beta-\gamma)\Pi(1-\gamma)}{\Pi(\alpha-\gamma)\Pi(\beta-\gamma)}
\]
Ebenso liefert die Subtraktion, wenn man die erste Gleichung mit \(\frac{\mathrm dP}{\mathrm dx}\) und die zweite mit \(P\) multipliziert,
\[
\frac{R\mathrm dP-P\mathrm dR}{\mathrm dx}=f(\alpha+1-\gamma,\beta+1-\gamma,2-\gamma)\frac{Q\mathrm dP-P\mathrm dQ}{\mathrm dx}
\]
und somit
\[
\begin{aligned}&C=(\gamma-1)f(\alpha+1-\gamma,\beta+1-\gamma,2-\gamma)\\&=\frac{\Pi(\alpha+\beta-\gamma)\Pi(\gamma-1)}{\Pi(\alpha-1)\Pi(\beta-1)}\end{aligned}
\]
Wenn man es vorzieht, lassen sich diese drei Gleichungen auch so darstellen, dass die Ableitungen \(\frac{\mathrm dP}{\mathrm dx},\ \frac{\mathrm dQ}{\mathrm dx},\ \frac{\mathrm dR}{\mathrm dx}\) durch endliche Funktionen ausgedrückt werden; so lautet zum Beispiel Formel 96
\SourceLeaf{232}
\[
\begin{aligned}
{}[97]\quad&\frac1\gamma F(\alpha,\beta,\alpha+\beta+1-\gamma,1-x)\\&\qquad F(\alpha+1,\beta+1,\gamma+1,x)\\
&+\frac1{\alpha+\beta+1-\gamma}F(\alpha+1,\beta+1,\alpha+\beta+2-\gamma,1-x)\\&\qquad F(\alpha,\beta,\gamma,x)\\
&=\frac{\Pi(\alpha+\beta-\gamma)\Pi(\gamma-1)}{\Pi\alpha\Pi\beta}x^{-\gamma}(1-x)^{\gamma-\alpha-\beta-1}
\end{aligned}
\]
\SixthArticle{52}
Bezeichnet man die Funktion $F(-\alpha,-\beta,1-\gamma,x)$ mit $S$, so ergibt sich
\[
0=\alpha\beta S-(1-\gamma+(\alpha+\beta-1)x)\frac{\mathrm dS}{\mathrm dx}-(x-xx)\frac{\mathrm{dd}S}{\mathrm dx^2}
\]
Verbindet man diese Gleichung mit I des vorigen Artikels, so ergibt sich
\[
0=\alpha\beta\left(\frac{S\mathrm dP+P\mathrm dS}{\mathrm dx}\right)
-(1-2x)\frac{\mathrm dP}{\mathrm dx}\mathbin{.}\frac{\mathrm dS}{\mathrm dx}
-(x-xx)\frac{\mathrm dS\mathrm{dd}P+\mathrm dP\mathrm{dd}S}{\mathrm dx^2}
\]
Diese Gleichung ist integrierbar und liefert
\[
\mathrm{Const.}=\alpha\beta PS-(x-xx)\frac{\mathrm dP}{\mathrm dx}\mathbin{.}\frac{\mathrm dS}{\mathrm dx}
\]
Der Wert der konstanten Größe $=\alpha\beta$, wie sich aus $x=0$ von selbst ergibt. Bevorzugen Sie eine endliche Form, so haben Sie
\[
\begin{aligned}
{}[98]\quad&F(\alpha,\beta,\gamma,x)F(-\alpha,-\beta,1-\gamma,x)\\
&-\frac{\alpha\beta}{\gamma-\gamma\gamma}\\&\qquad(x-xx)F(\alpha+1,\beta+1,\gamma+1,x)\\&\qquad F(1-\alpha,1-\beta,2-\gamma,x)=1
\end{aligned}
\]
Formen Sie hier jede der vier Funktionen nach Formel 82 um und schreiben Sie dann $\gamma-\alpha,\gamma-\beta$ für $\alpha,\beta$, so erhalten Sie
\[
\begin{aligned}
{}[99]\quad&(1-x)F(\alpha,\beta,\gamma,x)F(1-\alpha,1-\beta,1-\gamma,x)\\
&-\frac{\gamma-\alpha\mathbin{.}\gamma-\beta}{\gamma-\gamma\gamma}\\&\qquad xF(\alpha,\beta,\gamma+1,x)F(1-\alpha,1-\beta,2-\gamma,x)=1
\end{aligned}
\]
